\section{Discussion and Research Priorities}
\subsection{What action assurance adds}
The selected literature demonstrates increasingly capable interaction with tools and complex environments. It also makes the gap between a local success and a completed episode visible. Action assurance offers a common language for that gap. By separating proposal, admissibility, execution, and effect verification, CAA makes obligations testable without asserting that one architecture is always preferred. A clinician can reject a contract, a destination can reject a stale version, and a receipt can reveal a mismatch independently of how persuasive the planner's explanation appears.

This approach connects rather than replaces clinical reasoning and governance. Technical scope controls can prevent prohibited operations while still allowing a clinically wrong permitted proposal. Clinical review can improve judgment while still missing a stale resource or duplicate write. Institutions need both suitable reasoning evidence and enforceable effect boundaries. Administrative tasks likewise warrant attention: incorrect documentation or payer submissions can delay care even without direct treatment decisions.

\subsection{Tradeoffs and implementation gaps}
Fine-grained approval and verification add latency and clinician work. If every action is deferred, apparent safety may reflect an unusable system. If review is compressed into a generic approval button, its protection may be illusory. Evaluation should therefore measure review quality, time, interrupted work, delayed care, and responsibility for unresolved episodes. Optimal boundaries may differ across reversible administrative tasks, medication-related actions, and longitudinal management.

Interoperability helps expose data and operations but does not eliminate heterogeneous semantics. A standard resource may represent a request, an accepted task, or a completed effect depending on destination behavior. Cross-system preconditions and idempotency may be unavailable. Likewise, authenticated identity is necessary but not sufficient when upstream patient matching or encounter context is wrong. CAA should be adapted to a documented threat model and workflow, with residual uncertainty acknowledged.

\subsection{Limitations of this review}
Purposive selection favors reports that illuminate observable agency and execution; it cannot estimate the proportion of systems that meet any assurance requirement. Only PubMed and arXiv support the reproducible discovery inventory, with targeted retrieval adding primary conference and journal material. Terminology and indexing may omit relevant systems, and inclusion of preprints makes some evidence provisional. Agent versions evolve rapidly, so reported rates apply to cited configurations and snapshots.

AI-assisted extraction can miss methodological details and cannot stand in for independent clinical verification. One selected source is abstract-only. The interrupted original workspace required a rebuild; the current audit distinguishes refreshed records from unavailable historical exports. We do not provide formal risk-of-bias assessment, pooled efficacy estimates, or new clinical experiments. CAA synthesizes existing ideas and may overlap with other governance frameworks; its novelty lies in the explicit action/effect/authority lens and reviewable operational contract, not a claim of first invention. Its synthetic implementation is deliberately incomplete for production.

\subsection{Outstanding questions}
\begin{enumerate}
\item Which decision-time signals predict unsafe action effects and clinically important omissions, rather than only diagnostic disagreement?
\item How can approvals remain valid across changing evidence and multiple clinical systems with limited transactional support?
\item What level of human review improves safety without unacceptable interruption, delay, or automation bias?
\item How should memory, model, retrieval, and policy changes trigger revalidation, including subgroup coverage and access?
\item Which prospective designs can identify benefit and harm attributable to a specific bounded authority, with adequate follow-up and accountable recovery?
\end{enumerate}
